\chapter{Sukshma and Sthula Vyayama}
\label{ch:vyayama}

\emph{How does a class warm up before the asanas, and what is the difference between the gentle and the strong loosening practices?}
The syllabus names two families of loosening practice. Sukshma vyayama is a series of small, slow joint movements; sthula vyayama is a pair of larger, stronger movements for the whole body. A candidate performs both and knows their Sanskrit names. The idea behind them is introduced with the theory of vyayama and surya namaskara (see \cref{ch:vyayama-intro}).

\begin{unverified}
The steps of the Board's named practices follow an edited reprint of Swami Dhirendra Brahmachari's \emph{Yogic Suksma Vyayama and Yogic Sthula Vyayama} (the Board names the originals); the neck, trunk-twist and knee practices follow the Common Yoga Protocol. The reprint is still to be set against the original or the MDNIY monograph, and no text gives cautions for most of the named practices. A certified teacher must review this chapter before print.
\end{unverified}

\section{Sukshma vyayama}

\subsection*{Meaning and purpose}
Sukshma vyayama are loosening practices, also called micro-circulation practices: small, slow movements of one joint at a time, done standing or sitting. The standing ones start from samasthiti, the alert standing posture, which is the same as tadasana. Each practice in the Board's list is called a \emph{shakti vikasaka}, a practice that develops strength in one part. They make the joints mobile and prepare the body for stronger work, which is why they open a session. They are traditionally credited with joint mobility.

The syllabus list runs from the head down: the neck (griva shakti vikasaka), the shoulders (bhuja valli and purna bhuja shakti vikasaka), the trunk (kati shakti vikasaka), the thighs (jangha shakti vikasaka), the knees (janu shakti vikasaka), and the ankles and feet (pada-mula shakti vikasaka, and the gulpha, pada-prishtha and pada-tala practice).

\subsection*{How to teach them}
Teach every movement in two parts: first the position, then the breath that goes with it. Keep the range small and stop at pain. Watch for tension in the shoulders and the jaw. Offer a chair version to older students and to anyone who cannot stand comfortably.

\subsection*{Neck: griva shakti vikasaka}
The common protocol gives four stages, two rounds each. The rule for the breath is simple: breathe out as the head goes down or turns, breathe in as it returns.
\begin{enumerate}
\item \textbf{Forward and backward bending.} Breathe out and bring the chin towards the chest; breathe in and take the head back as far as is comfortable.
\item \textbf{Side bending.} Breathe out and take the ear towards the shoulder without lifting the shoulder; breathe in and come up; then the other side.
\item \textbf{Twisting.} With the head upright, breathe out and turn the chin towards one shoulder; breathe in and come back; then the other side.
\item \textbf{Rotation.} Breathe out and drop the head forward, then draw one slow full circle, breathing in as the head goes back and out as it comes forward; then circle the other way.
\end{enumerate}
Keep the shoulders relaxed and still throughout.

\begin{caution}
Elderly people and people with cervical problems or high blood pressure should avoid the full neck practice. People with neck pain may do it gently, seated, in a small range.
\end{caution}

\ycbfig{ycb1-fig-greeva}
\ycbfig{ycb1-fig-greevabend}

\subsection*{Shoulders: bhuja valli and purna bhuja shakti vikasaka}
\textbf{Bhuja valli.} Feet together. Raise one arm sideways above the head, palm open and facing out, and bring it down without touching the head or the thigh. Right arm, left arm, then both.

\textbf{Purna bhuja.} Make fists with the thumbs tucked in. Swing one arm in a full circle forward, then backward; then the other arm; then both arms. (The reprint holds the breath here; for Level 1 leave the breath natural until the original is read.)

\ycbfig{ycb1-fig-bhuja}

\subsection*{Trunk: kati shakti vikasaka}
The reprint gives five forms:
\begin{enumerate}
\item One fist behind the back, the other hand holding that wrist. Bend back on the in-breath; bend forward on the out-breath towards the knees.
\item Legs wide apart, hands on the hips. Bend back, then forward towards the floor.
\item Feet together. Quick bends back and forward (the reprint says with a jerk; no source settles whether this suits people with back trouble, so go slowly).
\item Arms stretched out to the sides. Bend the trunk to the left and to the right towards the calf, without moving the arms up or down. This is a side bend, not a twist.
\item Feet about two feet apart. Swing the trunk and the outstretched arms in a half circle to each side.
\end{enumerate}
The common protocol has its own trunk twist, which is not one of these five: feet two to three feet apart, arms at chest height with the palms facing; breathe out and twist to the left so that the right palm reaches the left shoulder, breathe in and come back, then the right; three rounds.

\begin{caution}
For the protocol twist: heart patients take care; avoid it with severe back pain, spine or disc disorders, after abdominal surgery and during menstruation. No source gives cautions for the five reprint forms; teach them slowly and stop at pain.
\end{caution}

\ycbfig{ycb1-fig-kati}

\subsection*{Thighs: jangha shakti vikasaka}
The syllabus lists two forms, II-A and II-B:
\begin{itemize}
\item \textbf{II-A.} Feet together, arms straight out in front, parallel to the floor. Bend the knees, keeping them together, until the thighs are level; heels and toes stay down. Rise slowly.
\item \textbf{II-B.} Heels raised, arms out to the sides. Bend with the knees apart, without sitting on the heels. Rise slowly.
\end{itemize}
Jumping jacks are another form in the reprint, but they are not in the syllabus list; leave them out.

\subsection*{Knees: janu shakti vikasaka}
Standing, swing one foot forward with a jerk of the knee, then back so that the heel touches the buttock; then the other leg. Ten times to begin. The common protocol's knee movement is a slow squat: breathe in and raise the arms to shoulder height, palms down; breathe out and bend the knees until the arms and thighs are parallel to the floor; breathe in and straighten; breathe out and lower the arms; three rounds.

\begin{caution}
Avoid the protocol squat in acute arthritis.
\end{caution}

\ycbfig{ycb1-fig-jangha}
\ycbfig{ycb1-fig-jangha-1}
\ycbfig{ycb1-fig-jangha-2}
\ycbfig{ycb1-fig-janu-1}
\ycbfig{ycb1-fig-janu-2}

\subsection*{Ankles and feet}
\textbf{Pada-mula A.} On the toes, rise and lower in a spring-like movement. \textbf{Pada-mula B.} Balanced on the toes, jump up and land on the toes.
\textbf{Gulpha-pada-prishtha-pada-tala.} Hold one foot about nine inches off the ground and circle the ankle each way; then the other foot.

\ycbfig{ycb1-fig-pada}
\ycbfig{ycb1-fig-gulpha}

\subsection*{The movements behind the practices}
The workshop shows the four kinds of movement that sukshma vyayama works through: flexion and extension (bending and straightening), abduction and adduction (moving away from and towards the midline), and rotation. Naming the movement helps a teacher explain what each practice does.

\ycbfig{ycb1-fig-flexors}
\ycbfig{ycb1-fig-extensors}
\ycbfig{ycb1-fig-abductors}
\ycbfig{ycb1-fig-adductors}
\ycbfig{ycb1-fig-rotators}
\ycbfig{ycb1-fig-sukshma}

\section{Sthula vyayama}

\subsection*{Meaning and purpose}
\emph{Sthula} means gross, and the reprint renders its own title as ``yogic gross exercise''. Where sukshma vyayama is gentle and works the joints one by one, sthula vyayama is a stronger loosening exercise for the whole body. The syllabus names two sthula practices: sarvanga pushti and hrid gati. Ordinary gym exercises such as push-ups, squats, planks, lunges and leg raises are sometimes listed under this heading; they are not in the Board syllabus.

\subsection*{Sarvanga pushti}
The name means ``whole-body nourishing''.
\begin{enumerate}
\item Stand with the feet as wide apart as possible. Make fists with the thumbs tucked in and place one wrist on the other.
\item Bend towards the right ankle.
\item Swing the joined arms up, back and round in a circle with the trunk, and bring them down towards the left ankle.
\item Move very slowly, then repeat in the other direction.
\end{enumerate}
The reprint's breath cue contradicts itself, so no breath cue is given here; no source gives the number of rounds.

\ycbfig{ycb1-fig-sthula-sarvanga}

\subsection*{Hrid gati (engine daud)}
\begin{enumerate}
\item Stand with the feet together, the elbows bent, fists with the thumbs tucked in.
\item Run about fifty small steps forward on the toes, the heels flicking back towards the buttocks, the arms thrusting forward and back alternately like pistons.
\item Run the same distance backward.
\item Breathe sharply and deeply through the nose throughout; the breath gives an engine-like hiss, hence the name engine daud, the locomotive exercise.
\end{enumerate}

\ycbfig{ycb1-fig-sthula-hrid}
\ycbfig{ycb1-fig-sthula}

\begin{caution}
No text read gives a caution list for either sthula practice. As general advice, not from a text: with high blood pressure, a heart condition, or back or neck pain, take advice from a qualified teacher or doctor first.
\end{caution}

\begin{recall}
\begin{itemize}
\item Sukshma: small, slow joint movements; sthula: stronger whole-body movements.
\item Neck: four stages, two rounds; breathe out as the head goes down or turns.
\item Series from the head down: griva, bhuja valli and purna bhuja, kati, jangha, janu, pada-mula and gulpha.
\item Bhuja valli raises the arm sideways; purna bhuja swings it in a full circle.
\item Sthula practices: sarvanga pushti and hrid gati (engine daud).
\end{itemize}
\end{recall}

\subsection*{Worked examples}
\begin{itemize}
\item \textbf{An older student with a stiff neck joins a class.} Teach the four neck stages seated on a chair, in a small range with the shoulders relaxed, and ask them to stop at any pain.
\item \textbf{A student with high blood pressure wants to join hrid gati.} Leave them out of the sthula round; give them the seated sukshma practices instead and suggest they ask their doctor.
\item \textbf{Model answer: the difference between sukshma and sthula vyayama.} Sukshma vyayama is gentle and works the joints one at a time for mobility; sthula vyayama is a stronger exercise for the whole body. Both are done before the asanas.
\end{itemize}

\subsection*{Check yourself}
\begin{enumerate}
\item What are sukshma vyayama?
\item Give the breath rule for the neck practice.
\item Name one group who should avoid the neck practice.
\item How do bhuja valli and purna bhuja differ?
\item Which forms of jangha shakti vikasaka does the syllabus list?
\item Which joint does pada-mula shakti vikasaka work?
\item Name the two sthula vyayama practices of the syllabus.
\item How are the hands held in sarvanga pushti?
\item Why is hrid gati called engine daud?
\end{enumerate}

\subsection*{Answers}
\begin{enumerate}
\item Subtle joint-loosening practices that warm the body and prepare it for asanas.
\item Breathe out as the head goes down or turns; breathe in as it returns.
\item Elderly people, people with cervical problems, or people with high blood pressure.
\item Bhuja valli raises one arm sideways above the head and lowers it; purna bhuja swings the arm, in a fist, in a full circle.
\item II-A and II-B.
\item The ankle and foot (rising on the toes, form A; a jump on the toes, form B).
\item Sarvanga pushti and hrid gati (engine daud).
\item Fists with the thumbs tucked in, one wrist on the other.
\item The sharp, deep nasal breathing gives an engine-like hiss while the arms move like pistons.
\end{enumerate}

\paragraph{Sources.} Swami Dhirendra Brahmachari, \emph{Yogic Suksma Vyayama and Yogic Sthula Vyayama}, edited reprint (to be set against the original or the MDNIY monograph); Common Yoga Protocol, Ministry of AYUSH, 2nd revised edition, 2016 (neck practice, trunk twist, knee squat; it names Brahmachari's work as its reference for the loosening practices); the Board syllabus for Level 1 (the shakti vikasaka list and the two sthula practices).
